\section{Related Work}

Our work bridges three research communities---scaling laws, data curation practices, and data-centric AI---by providing a quantitative framework for data quality that was missing from all three.

\subsection{Scaling Laws for Foundation Models}

Neural scaling laws predict model performance from compute budget and architectural choices. Kaplan et al.~\cite{kaplan2020scaling} showed power-law relationships between validation loss and model size N, dataset size D, and compute C, enabling compute-optimal model design. Hoffmann et al.~\cite{hoffmann2022training} refined this with the Chinchilla scaling laws, finding that prior models were undertrained---the compute-optimal ratio is approximately $N^{0.5}$ tokens per parameter, not the $N^{0.1}$ assumed by GPT-3. This led to smaller, better-trained models like Chinchilla (70B parameters, 1.4T tokens) outperforming Gopher (280B parameters, 300B tokens) at lower inference cost.

However, both Kaplan and Chinchilla formulations assume uniform data quality: the loss function $L(N, D, C)$ treats dataset size D as token count without quality weighting. One trillion deduplicated tokens and one trillion with 40\% duplicates receive identical treatment. While Hoffmann et al. acknowledge that ``data quality matters'' in discussion, they provide no mechanism to incorporate it into the scaling law. Subsequent work~\cite{muennighoff2023scaling,bi2024deepseek} has explored scaling laws for multilingual and code-pretrained models, but continues to treat quality as exogenous.

\textbf{Our contribution:} We propose Q(D) as a measurable fourth dimension, demonstrating that quality explains 61\% of information density variance orthogonally to dataset size. This extends $L(N, D, C)$ toward $L(N, D, Q(D), C)$, where quality is no longer implicit but quantified.

\subsection{Data Curation in Foundation Model Pretraining}

Practitioner knowledge about data quality is extensive but heuristic. GPT-3~\cite{brown2020language} employed fuzzy deduplication to remove near-duplicate documents, observing qualitative improvements in sample efficiency. The Pile~\cite{gao2020pile} assembled 22 curated data sources emphasizing domain diversity, arguing that breadth of coverage improves downstream task generalization. LLaMA~\cite{touvron2023llama} and LLaMA-2 combined aggressive quality filters (perplexity-based document scoring, toxic content removal) with careful corpus mixing ratios determined through manual tuning and validation loss monitoring.

More recent work has systematized some curation decisions. RedPajama~\cite{computer2023redpajama} reproduced LLaMA's filtering pipeline at scale, documenting heuristics like ``remove documents with $>30\%$ duplicate n-grams'' and ``sample domains to match LLaMA's target distribution.'' Raffel et al.~\cite{raffel2020exploring} introduced C4 (Colossal Clean Crawled Corpus) with rule-based filters (language detection, profanity removal, boilerplate suppression) that became widely adopted for T5, UL2, and Flan-T5 pretraining.

Despite this wealth of empirical practice, \textbf{no prior work has validated which quality dimensions matter most via controlled experiments}. Is deduplication more important than diversity? Does perplexity filtering provide value beyond removing non-English text? Practitioners lack quantitative evidence to prioritize curation investments.

\textbf{Our contribution:} We isolate and measure four quality dimensions on controlled C4 subsets, demonstrating that deduplication is the strongest predictor ($r = 0.72$) while diversity ($r = 0.65$), perplexity ($r = 0.58$), and efficiency ($r = 0.53$) all contribute non-redundantly. This provides evidence-based priority ranking for resource allocation.

\subsection{Data-Centric AI and Quality Metrics}

The data-centric AI movement~\cite{ng2021data,mazumder2022data} argues that improving data often yields larger performance gains than architectural innovations, especially for production systems. Northcutt et al.~\cite{northcutt2021pervasive} showed that label noise significantly degrades model performance and proposed automated cleaning methods (confident learning). Sorscher et al.~\cite{sorscher2022beyond} demonstrated that careful data pruning---removing low-value examples---can match full-dataset performance with 50\% fewer training samples. Abbas et al.~\cite{abbas2023semdedup} extended this to language models, showing 40\% data reduction via perplexity-based pruning with minimal loss degradation.

However, this work remains \textbf{qualitative about ``quality'' itself}. Ng's ``quality over quantity'' framing provides intuition but no operational metric. Data pruning methods (Sorscher, Abbas) define quality as ``samples the current model finds easy/hard'' (perplexity-based), which is task-specific and requires a trained model---not suitable for pretraining from scratch. Confident learning targets label noise in supervised settings, inapplicable to unsupervised pretraining.

\textbf{Our contribution:} We define quality via information density---an intrinsic, task-agnostic property measurable before training begins. Our Q(D) metric combines deduplication (redundancy), diversity (coverage), perplexity (noise), and efficiency (signal density) into a composite score validated against information-theoretic ground truth (token entropy, n-gram redundancy, semantic diversity). Unlike pruning methods that require a trained model, Q(D) can be computed on raw text.

\subsection{Information Theory and Learning Efficiency}

Our work connects to information-theoretic perspectives on learning efficiency. Shannon entropy quantifies information content in data~\cite{shannon1948mathematical}, and rate-distortion theory~\cite{cover2006elements} predicts that compression (removing redundancy) improves signal transmission efficiency. In the context of neural networks, Tishby \& Zaslavsky~\cite{tishby2015deep} argued that learning can be understood as information compression, and Fisher information~\cite{ly2017tutorial} measures the amount of information an observable random variable carries about an unknown parameter.

Recent work has applied information theory to dataset quality. Ash et al.~\cite{ash2020deep} used gradient-based scores to select informative training examples. Mindermann et al.~\cite{mindermann2022prioritized} analyzed loss-based prioritized sampling for improving sample efficiency. However, these methods measure ``informativeness to a specific model'' rather than intrinsic data quality.

\textbf{Our contribution:} We measure information density via entropy and redundancy \textit{before} any model training, providing a model-agnostic quality metric. While our initial hypothesis included Fisher information as a gradient-based density proxy, we found it unstable in preliminary experiments (Section 4.2), leading us to rely on entropy-based measures.

\subsection{Positioning Summary}

Our work occupies the intersection of three research areas: we provide the \textbf{quantitative quality metric} that scaling laws lack, the \textbf{controlled experimental validation} that practitioner heuristics lack, and the \textbf{task-agnostic measurement framework} that data-centric AI lacks. The closest prior work is Chinchilla (scaling law rigor) and RedPajama (curation documentation), but neither quantifies quality's effect on learning efficiency. We are the first to validate a multi-component Q(D) metric against information density through controlled experiments suitable for scaling law integration.
